\section{Continuous Assessment 1 (CA-1) --- Set B}
\label{sec:ca1_set_b}

\LPUPaperHeader{Continuous Assessment (CA-1)}{B}{60 Minutes}{30 Marks}{CO1 (Variables \& Types), CO2 (Conditionals \& Functions)}

\begin{center}
\sffamily\textbf{\color{AlertRed}Note: Attempt any five out of the six questions. Each question carries 6 marks.}
\end{center}

\questionitem{1}{
    Explain the LEGB (Local, Enclosing, Global, Built-in) scope resolution rule. Provide code examples demonstrating the use and distinct behaviors of the \texttt{global} and \texttt{nonlocal} keywords.
}{6}{CO2}{L2 Understand}

\questionitem{2}{
    Discuss floating-point representation artifacts in Python IEEE 754 arithmetic (e.g., \texttt{0.1 + 0.2 != 0.3}). Explain how to correctly perform financial arithmetic using the \texttt{decimal} module and tolerance comparisons via \texttt{math.isclose()}.
}{6}{CO1}{L3 Apply}

\questionitem{3}{
    Write a Python function \texttt{collatz\_sequence(n)} that generates the sequence of numbers for starting integer $n$ until it reaches 1 ($n \to n/2$ if even, $n \to 3n + 1$ if odd). Compute and return the peak value reached and total stopping time.
}{6}{CO2}{L3 Apply}

\questionitem{4}{
    Analyze Python's short-circuit evaluation mechanism for boolean operators \texttt{and} and \texttt{or}. What is printed by \texttt{print([] or 0 or "Hello" and 5)}? Explain the evaluation sequence step-by-step.
}{6}{CO1}{L4 Analyze}

\questionitem{5}{
    Explain the function of \texttt{random.seed()}. Differentiate between pseudo-random numbers and cryptographically secure random numbers, contrasting \texttt{random.choice()} with \texttt{secrets.choice()}.
}{6}{CO2}{L2 Understand}

\questionitem{6}{
    Write a recursive function to solve the classical \textbf{Tower of Hanoi} puzzle for $n$ disks. Draw the recursion tree for $n = 3$, showing disk movements and deriving total move count $M(n) = 2^n - 1$.
}{6}{CO2}{L4 Analyze}

\newpage
\subsection*{Solutions and Marking Rubrics --- CA-1 Set B}

\begin{solutionbox}[Solution to Question 1]
When resolving an identifier, Python searches namespaces in strict order:
\begin{enumerate}[leftmargin=1.5em]
    \item \textbf{Local}: Names assigned inside the current function body.
    \item \textbf{Enclosing}: Names in enclosing outer functions (closures).
    \item \textbf{Global}: Names defined at the top level of the current module file.
    \item \textbf{Built-in}: Preloaded symbols (\texttt{len}, \texttt{range}, \texttt{print}, \texttt{ValueError}).
\end{enumerate}
\textbf{Keyword Differences}:
\begin{pycode}
x = 10\\
def outer():\\
\_\_ESCTOKEN\_7\_\_ y = 20\\
\_\_ESCTOKEN\_7\_\_ def inner():\\
\_\_ESCTOKEN\_7\_\_ \_\_ESCTOKEN\_7\_\_ nonlocal y  \# Modifies enclosing outer's variable\\
\_\_ESCTOKEN\_7\_\_ \_\_ESCTOKEN\_7\_\_ global x    \# Modifies module-level global variable\\
\_\_ESCTOKEN\_7\_\_ \_\_ESCTOKEN\_7\_\_ y += 1; x += 1
\end{pycode}
\begin{markingrubric}
$\bullet$ Definition of LEGB lookup hierarchy: \textbf{2 Marks}\\
$\bullet$ Explanation and syntax of \texttt{global}: \textbf{2 Marks}\\
$\bullet$ Explanation and syntax of \texttt{nonlocal} in closures: \textbf{2 Marks}
\end{markingrubric}
\end{solutionbox}

\begin{solutionbox}[Solution to Question 2]
Floats are represented in binary scientific notation: $0.1_{10} = 0.0001100110011\dots_2$ which has an infinite repeating fractional component in base 2. Truncation to 53 bits results in a rounding error:
\begin{equation*}
0.1 + 0.2 = 0.3000000000000000444 \dots \neq 0.3
\end{equation*}
\textbf{Safe Comparison Solutions}:
\begin{pycode}
import math\\
from decimal import Decimal\\
~\\
\# Engineering tolerance check:\\
assert math.isclose(0.1 + 0.2, 0.3, rel\_tol=1e-9)\\
~\\
\# Exact financial calculation:\\
val = Decimal('0.1') + Decimal('0.2')  \# Decimal('0.3')
\end{pycode}
\begin{markingrubric}
$\bullet$ Binary representation explanation of repeating fraction: \textbf{2 Marks}\\
$\bullet$ Code implementation with \texttt{math.isclose}: \textbf{2 Marks}\\
$\bullet$ Exact monetary representation with \texttt{decimal.Decimal}: \textbf{2 Marks}
\end{markingrubric}
\end{solutionbox}

\begin{solutionbox}[Solution to Question 3]
\begin{pycode}
def collatz\_sequence(n):\\
\_\_ESCTOKEN\_7\_\_ if not (isinstance(n, int) and n >= 1):\\
\_\_ESCTOKEN\_7\_\_ \_\_ESCTOKEN\_7\_\_ raise ValueError("Input must be a positive integer >= 1")\\
\_\_ESCTOKEN\_7\_\_ seq = [n]\\
\_\_ESCTOKEN\_7\_\_ while n != 1:\\
\_\_ESCTOKEN\_7\_\_ \_\_ESCTOKEN\_7\_\_ if n \% 2 == 0:\\
\_\_ESCTOKEN\_7\_\_ \_\_ESCTOKEN\_7\_\_ \_\_ESCTOKEN\_7\_\_ n = n // 2\\
\_\_ESCTOKEN\_7\_\_ \_\_ESCTOKEN\_7\_\_ else:\\
\_\_ESCTOKEN\_7\_\_ \_\_ESCTOKEN\_7\_\_ \_\_ESCTOKEN\_7\_\_ n = 3 * n + 1\\
\_\_ESCTOKEN\_7\_\_ \_\_ESCTOKEN\_7\_\_ seq.append(n)\\
\_\_ESCTOKEN\_7\_\_ return \_\_ESCTOKEN\_5\_\_\\
\_\_ESCTOKEN\_7\_\_ \_\_ESCTOKEN\_7\_\_ "sequence": seq,\\
\_\_ESCTOKEN\_7\_\_ \_\_ESCTOKEN\_7\_\_ "stopping\_time": len(seq) - 1,\\
\_\_ESCTOKEN\_7\_\_ \_\_ESCTOKEN\_7\_\_ "peak\_value": max(seq)\\
\_\_ESCTOKEN\_7\_\_ \_\_ESCTOKEN\_6\_\_
\end{pycode}
\begin{markingrubric}
$\bullet$ Input validation and while loop termination condition: \textbf{2 Marks}\\
$\bullet$ Correct parity branching ($n//2$ vs $3n+1$): \textbf{2 Marks}\\
$\bullet$ Peak value and stopping time calculation: \textbf{2 Marks}
\end{markingrubric}
\end{solutionbox}

\begin{solutionbox}[Solution to Question 4]
Evaluation of \texttt{[] or 0 or "Hello" and 5}:
\begin{enumerate}[leftmargin=1.5em]
    \item Operator precedence: \texttt{and} (rank 12) binds tighter than \texttt{or} (rank 13). Expression parses as: \texttt{[] or 0 or ("Hello" and 5)}.
    \item Evaluate \texttt{"Hello" and 5}: \texttt{"Hello"} is non-empty (truthy). Python evaluates second operand: returns \texttt{5}.
    \item Evaluate \texttt{[] or 0 or 5}: Leftmost \texttt{[] or 0}: \texttt{[]} is empty list (falsy) $\to$ returns second operand \texttt{0}.
    \item Remaining: \texttt{0 or 5}: \texttt{0} is falsy $\to$ returns second operand \textbf{5}. Output is \textbf{5}.
\end{enumerate}
\begin{markingrubric}
$\bullet$ Operator precedence parsing (\texttt{and} over \texttt{or}): \textbf{2 Marks}\\
$\bullet$ Evaluation of \texttt{"Hello" and 5}: \textbf{2 Marks}\\
$\bullet$ Final left-to-right short-circuit step yielding 5: \textbf{2 Marks}
\end{markingrubric}
\end{solutionbox}

\begin{solutionbox}[Solution to Question 5]
\textbf{Seed Mechanics}: \texttt{random.seed(x)} initializes the internal state array (Mersenne Twister MT19937, state size 19937 bits). Given identical seeds, PRNG produces identical deterministic sequences across platforms.
\textbf{PRNG vs CSPRNG}:
\begin{itemize}[leftmargin=1.5em]
    \item \texttt{random}: Predictable after observing 624 consecutive 32-bit outputs. Never use for passwords, security tokens, or cryptographic keys.
    \item \texttt{secrets}: Uses OS entropy pools (\texttt{/dev/urandom} or Windows CryptGenRandom) ensuring non-reproducibility and resistance to state recovery.
\end{itemize}
\begin{markingrubric}
$\bullet$ PRNG state initialization and deterministic repeatability: \textbf{2 Marks}\\
$\bullet$ Vulnerability of MT19937 to state reconstruction: \textbf{2 Marks}\\
$\bullet$ Cryptographic security with OS entropy via \texttt{secrets}: \textbf{2 Marks}
\end{markingrubric}
\end{solutionbox}

\begin{solutionbox}[Solution to Question 6]
\begin{pycode}
def hanoi(n, source, auxiliary, target):\\
\_\_ESCTOKEN\_7\_\_ if n == 1:\\
\_\_ESCTOKEN\_7\_\_ \_\_ESCTOKEN\_7\_\_ print(f"Move disk 1 from \{source\} -> \{target\}")\\
\_\_ESCTOKEN\_7\_\_ \_\_ESCTOKEN\_7\_\_ return 1\\
\_\_ESCTOKEN\_7\_\_ moves = hanoi(n - 1, source, target, auxiliary)\\
\_\_ESCTOKEN\_7\_\_ print(f"Move disk \{n\} from \{source\} -> \{target\}")\\
\_\_ESCTOKEN\_7\_\_ moves += 1 + hanoi(n - 1, auxiliary, source, target)\\
\_\_ESCTOKEN\_7\_\_ return moves
\end{pycode}
\textbf{Recurrence Relation}: $M(n) = 2M(n-1) + 1$ with $M(1) = 1$.
Solving by substitution: $M(n) = 2^n - 1$. For $n = 3$, $M(3) = 2^3 - 1 = \mathbf{7\text{ moves}}$.
\begin{markingrubric}
$\bullet$ Recursive function definition with 3 steps: \textbf{2 Marks}\\
$\bullet$ Recurrence relation derivation $M(n) = 2^n - 1$: \textbf{2 Marks}\\
$\bullet$ Correct sequence of 7 steps for $n=3$: \textbf{2 Marks}
\end{markingrubric}
\end{solutionbox}

\newpage
